\icode{N23E-MG-SI-L3-261019-001}
\nceeID[4]{2023湖北}{2}
2022年12月8日，地球恰好运行到火星和太阳之间，且三者几乎排成一条直线，此现象被称为“火星冲日”。火星和地球几乎在同一平面内沿同一方向绕太阳做圆周运动，火星与地球的公转轨道半径之比约为\(3:2\)，如图所示。根据以上信息可以得出\paren[]

\begin{tikzpicture}[scale=.8,font=\scriptsize]
	\draw (0,0) circle(1.8);
	\draw (0,0) circle(2.5);
	\draw[dashed,cyan] (0,0)--(2.5,0);
	\draw[fill=orange!30] (0,0) circle(.15);
	\draw[fill=blue] (1.8,0) circle(.1);
	\draw[fill=red] (2.5,0) circle(.1);
	\draw[rotate=30,->,red] (2.7,0) arc(0:30:2.7);
	\draw[rotate=120,->,red](2,0) arc(0:40:2);
	\node at(-.15,0)[left]{太阳};
	\node at(1.7,.25)[left]{地球};
	\node at(2.6,0)[right]{火星};
\end{tikzpicture}
\begin{choices}
	\item 火星与地球绕太阳运动的周期之比约为\(27:8\)
	\item 当火星与地球相距最远时，两者的相对速度最大
	\item 火星与地球表面的自由落体加速度大小之比约为\(9:4\)
	\item 下一次“火星冲日”将出现在2023年12月8日之前
\end{choices}	
\begin{answer}
	B\qquad 提示：\(2\pi = (\frac{2\pi}{T_{earth}} - \frac{2\pi}{T_{mars}})t\)
\end{answer}